% Suggested response to reviewers and editor
% Manuscript: Brittle fracture in topology-optimized structures
% Submission ID: cc51f1c5-c888-45a4-8fae-ebefd5a3c8f7
% Review date: 18 August 2026
%
% This is a draft response and revision checklist. It intentionally does not
% modify the manuscript.

\documentclass[11pt,a4paper]{article}

\usepackage[T1]{fontenc}
\usepackage[utf8]{inputenc}
\usepackage{lmodern}
\usepackage{geometry}
\usepackage{microtype}
\usepackage{xcolor}
\usepackage{enumitem}
\usepackage{amsmath,amssymb}
\usepackage{hyperref}

\geometry{margin=25mm}
\hypersetup{colorlinks=true,linkcolor=blue!60!black,urlcolor=blue!60!black}

\newcommand{\Evar}{\ensuremath{\mathbf{E}_{\mathrm{var}}}}
\newcommand{\Emin}{\ensuremath{\mathbf{E}_{\mathrm{min}}}}
\newcommand{\Emax}{\ensuremath{\mathbf{E}_{\mathrm{max}}}}
\newcommand{\comment}[1]{\par\noindent\textbf{Reviewer comment.} #1\par}
\newcommand{\response}[1]{\par\noindent\textbf{Suggested response.} #1\par}
\newcommand{\change}[1]{\par\noindent\textbf{Suggested manuscript change.} #1\par}

\setlength{\parindent}{0pt}
\setlength{\parskip}{0.45\baselineskip}
\setlist[itemize]{leftmargin=*,topsep=0.2\baselineskip,itemsep=0.1\baselineskip}

\title{Suggested Response to Reviewers and Editor}
\author{Alexander Schlueter, Jan Pravda, Dustin Roman Jantos, Philipp Junker, Ralf Mueller}
\date{18 August 2026}

\begin{document}
\maketitle

\section*{Cover Letter to the Editor}

Dear Professor Schroeder,

We thank you and the reviewers for the careful assessment of our manuscript
``Brittle fracture in topology-optimized structures''. We have prepared a revised
manuscript that addresses the comments point by point. In particular, we have
clarified the scope of the contribution, moderated conclusions that could be
read as overstated, expanded the discussion of limitations of the employed
phase-field fracture model, and added details on the numerical implementation,
mesh resolution, porosity-to-fracture-resistance mapping, and interpretation of
the topology-optimization results.

The revision keeps the study as a sequential design-and-assessment
investigation: stiffness-based topology optimization is performed first, and
fracture is then assessed on the fixed optimized structures. We now state this
scope more explicitly and avoid implying that fracture resistance is optimized
or experimentally validated in the present work.

A detailed point-by-point response is provided below. We hope that the revised
manuscript now addresses the concerns raised during review and is suitable for
further consideration in \emph{Archive of Applied Mechanics}.

Sincerely,

Alexander Schlueter, on behalf of all authors

\section*{General Response}

We thank both reviewers for their thoughtful and constructive feedback. The
comments helped identify places where the manuscript needed a clearer scope,
more careful wording, and more transparent numerical details. The most important
planned revisions are:
\begin{itemize}
  \item narrow the title and contribution statement so the paper is clearly
  presented as a numerical post-optimization fracture assessment of
  stiffness-optimized graded porous structures;
  \item clarify the distinction from prior phase-field/topology-optimization
  studies, especially works that directly optimize fracture resistance;
  \item acknowledge limitations of the spectral split, mobility regularization,
  diffuse crack resolution, and porosity-to-\(G_c\) mapping;
  \item add mesh and implementation details for both topology optimization and
  fracture simulations;
  \item replace stress plots or add principal tensile stress/principal direction
  information, because these are more directly connected to the spectral split;
  \item temper claims about ``efficient use of material'' and explain
  alternative interpretations of distributed/fat damage zones;
  \item add quantitative highlights to the conclusion and correct the minor
  typographical and notation issues.
\end{itemize}

\section*{Response to the Editor}

\comment{Please ensure the results are accurately reported, any overstated
conclusions are rewritten and the limitations of the work fully explained.}

\response{We agree and have revised the manuscript with this in mind. The
revision now states more explicitly that the work is a numerical
post-optimization assessment, not a fracture-resistance optimization method and
not an experimentally validated material model. We have moderated statements
that could be interpreted as broad claims about material efficiency or general
fracture performance, and we have added a dedicated limitations discussion
covering the spectral energy split, the rate-dependent mobility term, the
diffuse nature of the predicted crack zones, the absence of experimental
validation for the porosity-to-\(G_c\) relation, and the restriction to one
benchmark geometry and loading configuration.}

\change{Add a short limitations paragraph at the end of the phase-field model
section and a second limitations paragraph in the conclusion. Also revise
sentences in the abstract, introduction, results discussion, and conclusion that
currently imply general material-efficiency conclusions.}

\section*{Reviewer 1}

\subsection*{General Scope and Contribution}

\comment{The manuscript would benefit from a clearer articulation of its
distinction from prior studies, e.g. 10.1016/j.cma.2018.10.010 and
10.1016/j.compstruc.2020.106205.}

\response{We agree that the distinction should be stated more sharply. The
suggested references address topology optimization and graded/infill concepts,
including approaches where fracture or related performance measures are part of
the optimization. Our contribution is different in scope: we do not propose a
new fracture-resistance optimization method. Instead, we take structures
obtained from thermodynamic topology optimization of porous media, where
stiffness is optimized through separate density and porosity fields, and then
ask how these fixed designs behave under a subsequent brittle phase-field
fracture assessment when the optimized porosity field is mapped to both
stiffness and local fracture resistance. This sequential assessment highlights
how a stiffness-driven porosity distribution changes crack initiation, crack
paths, peak load, and work up to peak load.}

\change{Revise the final paragraphs of the introduction to explicitly compare
the present sequential workflow with topology-optimization methods that include
fracture resistance, phase-field fracture, or stress constraints directly in the
optimization objective/constraints. Add the two suggested references if they are
not already cited at the relevant location.}

\comment{The mapping of porosity to fracture toughness appears to be based
solely on numerical investigations, without experimental validation or detailed
microstructural justification.}

\response{This is correct, and the manuscript should be clearer about this
limitation. The \(G_c(\phi_r)\) relation is a phenomenological numerical
mapping motivated by previous simulations on simplified porous media with
circular pores. It is not experimentally calibrated in the present work and is
not claimed to be universal. We use it only for the post-optimization fracture
assessment to study the consequences of assigning lower fracture resistance to
more porous regions.}

\change{In the subsection ``Fracture toughness as a function of porosity'', add
that the relation is numerically motivated, microstructure-specific, and not
experimentally validated here. Explain that different pore topologies, ligament
mechanisms, and porosity ranges may lead to different \(G_c\)-porosity trends.}

\comment{The title suggests a primary focus on brittle fracture, but the paper
devotes equal attention to topology optimization and fracture analysis. This may
mislead readers regarding the manuscript's main contribution.}

\response{We agree that the current title is broad. We propose to use a more
specific title that reflects the sequential design-and-assessment character of
the work.}

\change{Suggested title: ``Post-optimization brittle-fracture assessment of
topology-optimized graded porous structures''. Alternative shorter title:
``Brittle-fracture assessment of topology-optimized graded porous structures''.}

\subsection*{Phase-Field Fracture Model}

\comment{The total energy functional includes external work, which is
inconsistent with the original Griffith/Francfort-Marigo variational setting
for fracture. Larsen's two-potential approach resolves this only in a staggered
scheme, which is not clearly implemented or discussed here.}

\response{We thank the reviewer for pointing this out. Our implementation is a
rate-regularized phase-field fracture evolution under prescribed displacement
loading, solved monolithically in pseudo-time. The functional in the manuscript
was intended as the potential used to derive equilibrium equations under the
applied boundary conditions, while the crack evolution is governed by the
elastic/fracture terms plus the mobility regularization. We agree that the
wording can incorrectly suggest that we minimize a Griffith/Francfort-Marigo
functional including external work in the original global variational sense.
We will therefore revise the formulation and discussion to distinguish the
classical variational fracture setting from the incremental, rate-regularized
boundary-value problem solved here.}

\change{Rewrite the paragraph around the total potential in the phase-field
section. State explicitly that the external-work term enters the mechanical
equilibrium under prescribed boundary conditions, while the manuscript does not
claim existence of global minimizers for the original Francfort-Marigo problem
under Neumann loading. Add the reviewer-suggested references on this issue and
explain that no Larsen two-potential staggered scheme is used here.}

\comment{The statement that crack growth is driven by minimization of elastic
plus fracture energy contradicts Eq. (21) and adopts a Griffith statement
originally given for systems with no change in external energy.}

\response{We agree. The statement is imprecise for the boundary-value problem
as written and should be corrected.}

\change{Replace the sentence with a more careful formulation such as:
``In the present rate-regularized phase-field formulation, damage evolution is
driven by the local energetic competition between the degraded tensile elastic
energy and the regularized fracture energy, while the external loading supplies
work to the system through the imposed boundary conditions.''}

\comment{The spectral decomposition is known to fail in predicting traction-free
crack surfaces under shear deformations. The known drawbacks should be
acknowledged.}

\response{We agree. The spectral split is widely used and was chosen for its
standard ability to suppress crack growth under compression, but it has known
limitations in mixed-mode and shear-dominated situations. This matters for our
benchmark because several members are bending- and shear-influenced. We will
add an explicit limitation and avoid interpreting the final diffuse damage zones
as fully resolved traction-free cracks.}

\change{Add a paragraph after the spectral split equations acknowledging that
the Miehe spectral decomposition can lead to nonphysical residual tractions or
inaccurate mixed-mode behavior, especially under shear. Cite improved
alternatives, including representative crack elements or physically motivated
decompositions. Also refer back to this limitation when discussing Figs. 14--18.}

\comment{When using the spectral energy split, the first principal stress and
principal direction are of more interest than volumetric and deviatoric
stresses.}

\response{We agree. The present volumetric/deviatoric plots give general
orientation about tension/compression and shear-dominated regions, but they are
not the most direct diagnostic for the spectral split.}

\change{Replace the volumetric/deviatoric stress figures or add a new figure
showing the maximum principal tensile stress and corresponding principal
direction before fracture. Revise the discussion accordingly. If figure count
must be limited, move volumetric/deviatoric plots to supplementary material and
keep the principal-stress plot in the main manuscript.}

\comment{What portion of dissipated energy is contributed by the mobility
approach? Can this portion be physically motivated? Why is artificial damping
preferred over the staggered solution approach?}

\response{The mobility contribution is numerical/rate regularization rather
than a separately calibrated physical dissipation mechanism. It was introduced
to stabilize crack nucleation from initially undamaged material in the
monolithic solution. We agree that its energetic contribution should be reported
and its role should not be overinterpreted physically. The current manuscript
already defines \(D_s\), but it should quantify its magnitude relative to
\(\Pi_{\mathrm{frac}}\) and total supplied work.}

\change{Compute and report
\(D_s/(\Pi_{\mathrm{frac}}+D_s)\) and \(D_s/W\) at peak load and at the final
stored state for the cases shown in the energy-response figure. Add a short
paragraph stating that the mobility term is used for numerical robustness and
that its contribution is monitored to assess the degree of rate-regularization
influence. Explain that a staggered solution was not used because the present
implementation solves the coupled fields monolithically; note this as a
modeling/numerical choice rather than a claim of superiority.}

\comment{Fig. 18a demonstrates that the phase-field approach has issues
predicting component failure/vanishing reaction forces. The phase-field
evolution fails in bending dominated regions.}

\response{We agree that the final loss of load-carrying capacity is not fully
captured for all cases. For this reason, the manuscript uses the peak reaction
force and work up to peak load as comparison measures and does not interpret
the post-peak response as a reliable complete component-failure simulation.
This should be made more explicit.}

\change{In the discussion of the reaction-force plots, add that post-peak
behavior and full separation are affected by the diffuse crack approximation,
spectral split, mobility regularization, and nonlinear solver termination.
State that the comparison therefore focuses on pre-peak and peak-load measures,
not on complete structural separation or vanishing residual force.}

\comment{The reviewer is not convinced by the conclusion that widespread
damage hints at more efficient use of material. The thick crack zone could also
come from drawbacks of the spectral energy decomposition and mobility approach.}

\response{We agree that this conclusion was too strong. Distributed damage can
indicate that more regions approach criticality, but in the present model it can
also reflect the spectral split, the selected phase-field length scale, and the
mobility regularization. We will revise the wording to avoid claiming material
efficiency from the final damage pattern alone.}

\change{Replace the current sentence with: ``The \Evar cases show more
spatially distributed degradation than the homogeneous-porosity references.
This indicates that the optimized porosity field changes where critical states
are reached, but the width and extent of the diffuse damage zones must be
interpreted with caution because they can also be influenced by the spectral
split, mobility regularization, and chosen phase-field length scale.''}

\comment{Irreversibility is based on \(s=0\). Is any threshold used in the
implementation? Several points may never reach perfect \(s=0\).}

\response{The implementation should be described more precisely. In practice,
irreversibility is imposed using a numerical threshold for the phase field,
rather than relying on exact floating-point equality.}

\change{Verify the actual threshold in the code and state it explicitly, e.g.
``Nodes/quadrature points with \(s \le s_{\mathrm{irr}}\) are fixed at
\(s=0\) in subsequent increments, with \(s_{\mathrm{irr}}=\ldots\).'' If no
threshold is currently used, change the implementation in the revision stage or
state that no threshold was used and discuss the limitation.}

\comment{The manuscript claims that \(\kappa_s\) controls residual stiffness in
broken material, but this is the role of the undegraded compressive energy
\(\psi_e^-\).}

\response{We agree that the statement needs to be more precise. The parameter
\(\kappa_s\) controls only the residual part of the degraded tensile energy
contribution. The compressive contribution remains undegraded through
\(\psi_e^-\), which also affects the mechanical response in cracked regions.}

\change{Revise the sentence to: ``The parameter \(\kappa_s\) prevents complete
loss of the degraded tensile stiffness contribution, while the compressive
energy \(\psi_e^-\) remains undegraded in the spectral split.''}

\comment{The mapping from porosity to fracture toughness is based solely on
numerical studies. Experimental evidence and microstructural mechanisms are not
referenced. The porosity range may cross transitions depending on
microstructure topology.}

\response{We agree and will clarify that the chosen relation is not a general
law for porous materials. The relation depends on the assumed circular-pore
microstructure and on the numerical experiment used in the previous study.
Other microstructures may show different mechanisms, including transitions
between tension/compression-dominated and bending-dominated ligament response.}

\change{Add a paragraph in the \(G_c(\phi_r)\) subsection discussing possible
microstructural-mechanism changes and stating that the present mapping should
be regarded as a model assumption for this numerical study. Add suitable
references to experimental or micromechanical studies on porous/cellular
material fracture if available.}

\comment{Explicit crack resolution is not achieved due to fat phase-field
regions. Similar results could perhaps be achieved by a non-local damage model
at lower computational cost.}

\response{We agree that the final phase-field zones are diffuse and should not
be interpreted as explicitly resolved sharp cracks. The phase-field model is
used because it naturally handles nucleation and path changes in complex
optimized geometries, but the present results do not prove that a non-local
damage model would be insufficient.}

\change{Add this as a limitation: the present phase-field simulations represent
regularized damage/crack zones whose width depends on \(\beta_s\), and a
comparison with non-local damage or cohesive-zone formulations is outside the
scope but would be valuable future work.}

\subsection*{Topology Optimization}

\comment{Fig. 4: The schematic does not adequately represent the staggered
solution scheme for topology optimization.}

\response{We agree that the current schematic is too schematic and does not
show all relevant update loops.}

\change{Revise Fig. 4 to show initialization, mechanical equilibrium solve,
sensitivity evaluation, separate \(\chi\)- and \(\phi\)-evolution/update steps,
constraint/KKT update, convergence checks for objective and design variables,
and the loop to the next iteration.}

\comment{There is no comparison of constant-porosity optimization results with
established literature.}

\response{We agree that a short contextual comparison would help. Since our
constant-porosity cases reduce to a density-based topology optimization problem
with different homogeneous stiffness values, the resulting truss-like
double-cantilever structures can be compared qualitatively with established
compliance-minimizing topology-optimization benchmarks. We will add such a
comparison without claiming one-to-one quantitative identity, because exact
geometries, loads, regularization, and constraints differ between studies.}

\change{Add a paragraph after the constant-porosity results comparing the
observed truss-like layouts with literature benchmark behavior for
compliance-minimizing topology optimization and noting the limitations of this
comparison.}

\comment{The manuscript lacks mesh homogeneity, element size, and number of
elements, which are critical for assessing the accuracy, especially for thin
structures.}

\response{We agree and will add the missing details. The topology-optimization
mesh uses \(600\times100\) structured quadrilateral elements over a
\(6\,\mathrm{mm}\times1\,\mathrm{mm}\) design domain, corresponding to a
uniform element size of \(0.01\,\mathrm{mm}\) in both directions. The fracture
domains are re-meshed with first-order triangular elements; the manuscript
should report representative element sizes and element counts for each
\(\Omega_f\).}

\change{Add a table with, for each fracture case, the number of triangular
elements, minimum/maximum/mean element edge length or area, and the ratio of
element size to \(\beta_s\). Also state explicitly that the topology mesh is
homogeneous with \(h_x=h_y=0.01\,\mathrm{mm}\).}

\comment{Intermediate-density regions in Figs. 5 and 7 contradict the intended
separation of topology and stiffness optimization.}

\response{We understand the concern. The method uses a density field to define
material presence and a porosity field to distribute stiffness within material,
but numerical regularization and finite penalization do not produce a perfectly
binary density field everywhere. The fracture simulations therefore extract
\(\Omega_f\) using the threshold \(\chi\ge0.5\), after which \(\chi\) is set to
one and stiffness variation is governed by \(\phi\). We will clarify that the
separation is conceptual and algorithmic, not a claim that the optimized
density field is exactly binary before thresholding.}

\change{Add this clarification near the definition of \(\Omega_f\) and in the
discussion of Figs. 5 and 7. Also report the amount of intermediate-density
area, e.g. area fraction with \(0.1<\chi<0.9\), if available.}

\comment{The porosity field shows clusters without larger transitions. The
potential of the method is not well represented if it simplifies to two
materials. The interaction of porosity length scale and phase-field length
scale cannot be sufficiently studied when porosity only contains sharp
transitions.}

\response{We agree that much of the optimized porosity field is close to its
bounds for the chosen material interpolation and parameter set. The method
allows graded porosity fields, and larger \(\beta_\phi\) indeed smooths the
transitions, but the present benchmark and stiffness objective favor
near-extreme porosity values. Therefore, the interaction between
\(\beta_\phi\) and \(\beta_s\) is studied only for the transitions that occur
in this benchmark, not for a fully smooth continuously graded material field.
We will state this limitation explicitly.}

\change{Revise the \(\beta_\phi\) discussion to say that the investigated
optimized fields are largely clustered near porosity bounds, so conclusions
about porosity/fracture length-scale interaction are limited to these
optimization outcomes. Also temper the sentence about the internal fracture
length scale needing to be small compared with geometric features by adding
that it should also be interpreted relative to material-field transition
widths.}

\comment{For the spectral split, principal tensile stress and direction are
more relevant than volumetric or deviatoric stresses.}

\response{We agree and will address this together with the analogous
phase-field comment above.}

\change{Add or replace stress visualization with maximum principal tensile
stress and principal directions.}

\section*{Reviewer 2}

\comment{The title is too broad and should be specified.}

\response{We agree. As also noted by Reviewer 1, the title should better
reflect the sequential design-and-assessment scope.}

\change{Change the title to ``Post-optimization brittle-fracture assessment of
topology-optimized graded porous structures'' or ``Brittle-fracture assessment
of topology-optimized graded porous structures''.}

\comment{The more general part of the introduction could be condensed.}

\response{We agree. The introduction can be shortened before the more directly
related literature discussion. This will also help the paper emphasize its
specific contribution more clearly.}

\change{Condense the broad introductory material on porous materials and
general topology optimization, keeping only the context necessary to motivate
graded porosity, stiffness optimization, and post-optimization fracture
assessment.}

\comment{Why was fracture resistance not incorporated directly into the
optimization problem? Why is this more challenging?}

\response{This is an important point. We chose a sequential strategy because it
separates two questions: first, what structures are obtained when porous
topology optimization targets stiffness; second, how do those fixed structures
perform under brittle fracture assessment when the same porosity field affects
local fracture resistance. Incorporating fracture resistance directly into the
optimization is substantially more challenging because crack initiation and
propagation introduce path dependence, irreversibility, non-smooth or
history-dependent objectives, high computational cost for each design update,
and sensitivity-analysis difficulties for evolving crack topologies.}

\change{Add this rationale to the end of the workflow paragraph in
Section 2. Emphasize that direct fracture-resistance optimization is a natural
future extension rather than the objective of the present paper.}

\comment{The contour for \(\Emin\) with \(\rho=0.3\) and \(\Emax\) with
\(\rho=0.6\) are the same, whereas \(\Emax\) with \(\rho=0.3\) differs
significantly. Can this be explained in more detail?}

\response{Yes. The similarity follows from the mass constraint and homogeneous
porosity values. In \(\Emin\), the material has the maximum allowed physical
porosity, so only half of the local material density is present in regions with
\(\chi=1\). Thus, \(\rho=0.3\) with \(\Emin\) allows a similar occupied
geometric domain as \(\rho=0.6\) with \(\Emax\), where material is fully dense.
By contrast, \(\Emax\) with \(\rho=0.3\) has fully dense material but only half
the allowed total mass, so the optimizer produces a more delicate structure
with fewer/thinner members.}

\change{Add this explanation after Fig. 7 and/or Table 2, including the
relation between \(\phi_\mathrm{min}=0.5\), homogeneous porosity, and the mass
constraint.}

\comment{For the more widespread \(\Evar\) damage patterns, it would be
interesting to see how the pattern evolves. Where is first crack onset, and how
does propagation proceed within the structure?}

\response{We agree that the final damage field alone does not show the crack
sequence. We will add a short evolution description and, if space permits, an
additional figure with selected phase-field snapshots for representative cases.
For the present simulations, crack onset is expected at tensile fixed-boundary
regions or in the highly loaded inclined members depending on \(\rho\),
\(\beta_s\), and the local porosity field; the exact onset locations should be
reported from the stored simulation data.}

\change{Inspect the stored phase-field histories and add either a new figure
with three or four snapshots for representative \(\Evar\) cases, or add a
concise text paragraph identifying first crack onset and subsequent propagation
for the cases in Fig. 17/related final-pattern figures.}

\comment{In the conclusions, highlight the most important quantitative results
for the different cases considered, e.g. 36--64\% increase.}

\response{We agree and will add the key numerical ranges to the conclusion.
This will make the main findings easier to extract.}

\change{Add quantitative statements such as: for \(\rho=0.3\), \(\Evar\)
increases \(\max R_y\) by approximately 36--64\% relative to \(\Emin\) over
the considered \(\beta_s\) range; for \(\rho=0.6\), \(\Evar\) gives peak loads
comparable to or slightly above \(\Emax\); the stiffness gain of \(\Evar\)
relative to the homogeneous-porosity alternatives is about 2.5\% for
\(\rho=0.3\) and about 8\% for \(\rho=0.6\). Verify exact percentages before
final revision.}

\subsection*{Minor Corrections}

\comment{In some references, journals are not given, but an access date is
given. Add journal information and remove the date.}

\response{We will check and correct the bibliography entries.}

\change{Clean the BibTeX files by adding missing journal fields where
available and removing unnecessary access dates for journal articles.}

\comment{In equations (5)--(7), upper and lower indices seem mixed up.}

\response{We will correct the notation.}

\change{Check Eqs. (5)--(7) and make superscripts/subscripts consistent with
the surrounding definitions.}

\comment{In equations (5), (8), and (9), brackets seem to be missing.}

\response{We will correct the bracket notation.}

\change{Review Eqs. (5), (8), and (9) and add missing brackets/parentheses.}

\comment{It seems inconsistent to denote \(p\) as a penalization factor, as it
is an exponent rather than a factor.}

\response{We agree.}

\change{Replace ``penalization factor'' by ``penalization exponent'' wherever
\(p\) is discussed.}

\comment{For the parameters before Eq. (14), it does not make sense to give
decimal places for large numbers such as \(\mu_1\) and \(\lambda_1\).}

\response{We agree.}

\change{Round \(\mu_1\) and \(\lambda_1\) to sensible precision, e.g.
\(\mu_1\approx 8.08\cdot10^4\,\mathrm{N/mm^2}\) and
\(\lambda_1\approx 1.21\cdot10^5\,\mathrm{N/mm^2}\), or state only the
formulas and \(E_1,\nu\).}

\comment{In Eq. (24), why is the summation until \(n\), but not until 3? If it
is not 3, use another symbol because \(n\) denotes the number of finite
elements in \(x\)-direction.}

\response{We agree. The spectral decomposition should use the spatial
dimension as summation limit and avoid conflict with the mesh notation.}

\change{Replace \(\sum_{i=1}^{n}\) by \(\sum_{i=1}^{d}\), with \(d\) the
spatial dimension, or by \(\sum_{i=1}^{2}\) if the implementation and
presentation are strictly two-dimensional plane strain.}

\comment{There must not be a subsection 2.2.1 without a following subsection
2.2.2.}

\response{We agree.}

\change{Change ``2.2.1 Fracture toughness as a function of porosity'' to an
unnumbered paragraph heading or create a matching second subsubsection if the
section structure is expanded.}

\comment{On page 23, the sentence ``Recall that in all cases the same mass of
material is employed,'' should conclude with a period instead of a comma.}

\response{We agree and will correct the punctuation.}

\change{Replace the comma by a period.}

\section*{Revision Checklist}

\begin{itemize}
  \item Decide final title.
  \item Add/adjust citations requested by Reviewer 1, especially the two
  topology/fracture-related works and the phase-field variational/spectral-split
  limitation references.
  \item Rewrite phase-field potential discussion and energy-minimization
  wording.
  \item Add spectral-split limitations and replace/add stress diagnostics using
  principal tensile stress and direction.
  \item Quantify mobility dissipation contribution \(D_s\).
  \item Verify and report irreversibility threshold.
  \item Add fracture mesh statistics and topology mesh element size.
  \item Add explanation of \(\Emin,\rho=0.3\) versus \(\Emax,\rho=0.6\).
  \item Add crack-evolution snapshots or text from stored phase-field histories.
  \item Moderate claims about distributed damage/material efficiency.
  \item Add quantitative conclusions and all minor corrections.
\end{itemize}

\end{document}
